\chapter{Normalteiler, Homomorphismen und Quotienten}

\section{Normalteiler und Normalabschluss}

\begin{definition}[Normalteiler]
Eine Untergruppe \(N\leq G\) heißt \emph{normal}, geschrieben
\(N\trianglelefteq G\), wenn \(gNg^{-1}=N\) für jedes \(g\in G\) gilt.
\end{definition}

\begin{theorem}[Normalitätskriterien]
Für eine Untergruppe \(N\leq G\) sind äquivalent:
\begin{enumerate}
\item \(N\trianglelefteq G\).
\item \(gNg^{-1}\subseteq N\) für alle \(g\in G\).
\item \(gN=Ng\) für alle \(g\in G\).
\item Die Vorschrift \((gN)(hN)=ghN\) auf \(G/N\)
      ist unabhängig von der Wahl der Repräsentanten.
\item \(N\) ist der Kern eines Gruppenhomomorphismus mit Definitionsgruppe \(G\).
\end{enumerate}
\cite{milne,conradquot}
\end{theorem}

\begin{theorem}[Grundkonstruktionen mit Normalteilern]
Durchschnitte von Normalteilern von \(G\) sind normal in \(G\).
Für \(N\trianglelefteq G\) und \(H\leq G\) gelten
\[
 HN=NH\leq G,\qquad H\cap N\trianglelefteq H.
\]
Für \(N,M\trianglelefteq G\) ist außerdem \(NM\trianglelefteq G\).
Jede Untergruppe \(H\) mit \([G:H]=2\) ist normal, auch wenn
\(G\) unendlich ist. \cite{milne,conradquot}
\end{theorem}

\begin{definition}[Normalabschluss]
Für \(A\subseteq G\) heißt
\[
 \langle\!\langle A\rangle\!\rangle_G
 =\bigcap\{N\trianglelefteq G:A\subseteq N\}
\]
der \emph{Normalabschluss} von \(A\) in \(G\).
\end{definition}

\begin{theorem}[Erzeugung des Normalabschlusses]
Für jede Teilmenge \(A\subseteq G\) gilt
\[
 \langle\!\langle A\rangle\!\rangle_G
 =\langle gag^{-1}:a\in A,\ g\in G\rangle.
\]
Er ist der kleinste Normalteiler von \(G\), der \(A\) enthält.
Für \(A=\varnothing\) ist er \(\{1\}\). \cite{milne}
\end{theorem}

\section{Homomorphismen, Kerne und Bilder}



\begin{theorem}[Automatisch erhaltene Gruppendaten]
Für jeden Homomorphismus \(f:G\to H\), jedes \(g\in G\)
und jedes \(n\in\mathbb Z\) gelten
\[
 f(1_G)=1_H,\qquad f(g^{-1})=f(g)^{-1},\qquad f(g^n)=f(g)^n.
\]
Die Inverse eines Gruppenisomorphismus ist wieder ein
Gruppenhomomorphismus; Kompositionen von Gruppenhomomorphismen
sind Gruppenhomomorphismen. \cite{project40,milne}
\end{theorem}

\begin{theorem}[Untergruppenbilder und Normalteilerbilder]
Für einen Gruppenhomomorphismus \(f:G\to H\), \(U\leq G\) und \(V\leq H\) gelten
\[
 f(U)\leq H,\qquad f^{-1}(V)\leq G,\qquad
 \ker f\trianglelefteq G.
\]
Ist \(V\trianglelefteq H\), so ist \(f^{-1}(V)\trianglelefteq G\).
Ist \(U\trianglelefteq G\), so ist \(f(U)\trianglelefteq f(G)\).
Normalität von \(f(U)\) in ganz \(H\) folgt insbesondere dann,
wenn \(f(G)=H\). \cite{conradhom,milne}
\end{theorem}

\begin{theorem}[Kernkriterium und Fasern]
Für einen Gruppenhomomorphismus \(f:G\to H\) und \(x,y\in G\) gilt
\[
 f(x)=f(y)\iff x^{-1}y\in\ker f.
\]
Jede nichtleere Faser ist damit eine Nebenklasse:
\[
 f^{-1}(\{f(x)\})=x\ker f=(\ker f)x.
\]
Insbesondere ist \(f\) genau dann injektiv, wenn \(\ker f=\{1_G\}\).
\cite{conradhom}
\end{theorem}

\begin{theorem}[Bestimmtheit durch Erzeuger]
Ist \(A\subseteq G\) mit \(\langle A\rangle=G\), so sind zwei
Homomorphismen \(f,h:G\to K\) bereits gleich, sobald
\(f(a)=h(a)\) für alle \(a\in A\) gilt.
Für jeden Homomorphismus \(f:G\to H\) und jedes \(A\subseteq G\) gilt
\[
 f(\langle A\rangle)=\langle f(A)\rangle.
\]
\cite{milne,conradhom}
\end{theorem}

\section{Quotientengruppe und Isomorphiesätze}

\begin{definition}[Quotientengruppe und kanonische Projektion]
Für \(N\trianglelefteq G\) trägt \(G/N\) die Verknüpfung
\((gN)(hN)=ghN\). Die Abbildung
\[
 q_N:G\longrightarrow G/N,\qquad g\longmapsto gN
\]
heißt \emph{kanonische Projektion}.
\end{definition}

\begin{theorem}[Quotientensatz und universelle Eigenschaft]
\label{b40guide:thm:quotient}
Für \(N\trianglelefteq G\) ist \(G/N\) eine Gruppe mit Einheit \(N\)
und Inversen \((gN)^{-1}=g^{-1}N\).
Die Projektion \(q_N\) ist ein surjektiver Homomorphismus mit Kern \(N\).
Für jeden Homomorphismus \(f:G\to H\) sind äquivalent:
\[
 N\subseteq\ker f
 \quad\Longleftrightarrow\quad
 \exists!\,\overline f:G/N\to H
 \text{ Homomorphismus mit }f=\overline f\circ q_N.
\]
Dann gilt \(\overline f(gN)=f(g)\).
\cite{milne,conradquot}
\end{theorem}

\begin{theorem}[Homomorphiesatz; erster Isomorphiesatz]
\label{b40guide:thm:first-iso}
Für jeden Gruppenhomomorphismus \(f:G\to H\) ist
\[
 G/\ker f\longrightarrow\operatorname{im}f,\qquad
 g\ker f\longmapsto f(g)
\]
ein Isomorphismus. Insbesondere ist jedes homomorphe Bild von \(G\)
eine Quotientengruppe von \(G\), bis auf Isomorphie. \cite{milne}
\end{theorem}

\begin{theorem}[Zweiter Isomorphiesatz]
Für \(H\leq G\) und \(N\trianglelefteq G\) ist
\[
 H/(H\cap N)\longrightarrow HN/N,\qquad
 h(H\cap N)\longmapsto hN
\]
ein Isomorphismus. \cite{milne}
\end{theorem}

\begin{theorem}[Dritter Isomorphiesatz]
Für \(N\subseteq M\) mit \(N,M\trianglelefteq G\) gilt
\(M/N\trianglelefteq G/N\), und
\[
 (G/N)/(M/N)\longrightarrow G/M,\qquad
 (gN)(M/N)\longmapsto gM
\]
ist ein Isomorphismus. \cite{milne}
\end{theorem}

\begin{theorem}[Untergruppenkorrespondenz]
Für \(N\trianglelefteq G\) sind
\[
 U\longmapsto U/N,\qquad V\longmapsto q_N^{-1}(V)
\]
zueinander inverse Zuordnungen zwischen den Untergruppen
\(N\leq U\leq G\) und den Untergruppen \(V\leq G/N\).
Sie erhalten Inklusionen, Durchschnitte und erzeugte Untergruppen.
Außerdem gelten
\[
 U\trianglelefteq G\iff U/N\trianglelefteq G/N,\qquad
 [G:U]=[G/N:U/N].
\]
\cite{milne}
\end{theorem}

\section{Kommutatoren und der kanonische abelsche Quotient}

\begin{definition}[Kommutator und Kommutatoruntergruppe]
Für \(a,b\in G\) sei \([a,b]=aba^{-1}b^{-1}\).
Für Untergruppen \(A,B\leq G\) sei
\[
 [A,B]=\langle [a,b]:a\in A,\ b\in B\rangle.
\]
Insbesondere heißt \(G'=[G,G]\) die \emph{Kommutatoruntergruppe}
von \(G\). Die Gruppe \(G_{\mathrm{ab}}=G/G'\) heißt
\emph{Abelisierung} von \(G\).
\end{definition}

\begin{theorem}[Kommutatoren und Normalteiler]
Es gilt \([a,b]=1\iff ab=ba\).
Für \(A,B\trianglelefteq G\) gelten
\[
 [A,B]\trianglelefteq G,\qquad [A,B]\subseteq A\cap B.
\]
Die Untergruppe \(G'\) wird von jedem Automorphismus von \(G\)
auf sich abgebildet und ist insbesondere normal in \(G\).
\cite{milne}
\end{theorem}

\begin{theorem}[Charakterisierung aller abelschen Quotienten]
Für jeden Normalteiler \(N\trianglelefteq G\) gilt
\[
 G/N\text{ ist abelsch}\iff G'\subseteq N.
\]
Die Abelisierung \(G_{\mathrm{ab}}\) ist abelsch, und jeder
Homomorphismus \(f:G\to A\) in eine abelsche Gruppe faktorisiert
eindeutig durch \(G\to G_{\mathrm{ab}}\).
Diese Konstruktion ist für jede Ausgangsgruppe \(G\) definiert.
\cite{milne}
\end{theorem}

\begin{theorem}[Abelisierung und Homomorphismen]
Jeder Homomorphismus \(f:G\to H\) erfüllt \(f(G')\subseteq H'\)
und induziert einen Homomorphismus
\[
 f_{\mathrm{ab}}:G/G'\longrightarrow H/H',\qquad
 gG'\longmapsto f(g)H'.
\]
Bei \(f(G)=H\) gilt \(f(G')=H'\). \cite{milne}
\end{theorem}
